% ============================================================================
% BOLUM 6: SONUC
% ============================================================================
\section{Sonuç}
\label{sec:conclusion}

CNN--Görü Dönüştürücüsü gürbüzlük karşılaştırmalarındaki uyuşmazlığın ne kadarının modellerden değil ölçümden geldiğini sorduk ve bunu; modelleri, veriyi, tehdit modelini ve saldırı bütçesini sabit tutup yalnızca karşılaştırmanın nasıl puanlandığını değiştirerek yanıtladık.

Ölçüm hükmediyor. Özdeş çekişmeli eğitilmiş ResNet-18 ve ViT-Tiny çiftlerinde, mimariler arası transfer asimetrisi dört yerleşik koşullama protokolü boyunca $+4{,}4$ ile $+14{,}6$ puan arasında değişmektedir --- 3,3 katlık bir yayılım; aynı nicelik üzerinde ölçüldüğünde bu protokol etkisi, bağımsız eğitim koşuları arasındaki standart sapmanın $3{,}3$ ile $22{,}7$ katı arasındadır --- ve protokol, bir eşdeğerlik testinin kabul mü ret mi edileceğine karar vermektedir. Koşulsuz oranlar böyle bir karşılaştırma için hiçbir biçimde geçerli temel değildir: WideResNet-28-10 referansının eklendiği $3\times3$ matriste, koşullu oranlardan sapmaları hedefin temiz hatasıyla $r = 0{,}997$ ile açıklanmaktadır. Protokolün değiştirmediği şey yöndür; yön protokoller, tohumlar ve saldırı aileleri boyunca kararlıdır ve üç hedefli analizimiz onu CNN'de üretilen pertürbasyonların özel bir gücünden çok hedefin kendi kırılganlığıyla ($r = 0{,}986$) ilişkilendirmektedir.

Koşullu ayrıştırma, mutlak gürbüzlüğü temiz doğruluk çarpanı ile koşullu duyarlılık çarpanına ayırır. Düzeltilmiş koşularımızda ikisi de CNN lehinedir; ancak mutlak sıralaması özdeş olan önceki bir kontrol noktası kümesinde iki çarpan arasındaki pay farklıydı. Bu, mekanizma düzeyindeki iddiaların sıralama düzeyindeki iddialardan daha kırılgan olduğunu ve desteklenmek için daha çok eğitim koşusu gerektirdiğini gösterir.

Bu literatürde yaygın üç iddia doğrudan ölçümden sağ çıkamamıştır: çekişmeli eğitilmiş CNN gradyanları daha seyrektir ama mekânsal olarak daha lokalize değildir; CLS dikkat entropisi $n=1000$'de $L_\infty$ saldırısı altında değişmemektedir; dönüştürücüye özgü bir transfer saldırısı bu rejimde transfer kazancı sağlamamaktadır. İki davranışsal fark ise bağımsız kontrol noktası kümelerinde tekrarlanarak ayakta kalmıştır: mimarinin değil çekişmeli eğitimin yarattığı gradyan seyrekliği sıralaması ile çekişmeli hasarın derinlik profili --- CNN'de geç, yerel ve büyüklük değiştiren bir olay; ViT'te dağıtık, yalnızca yön değiştiren bir dönme.

Bu ölçümlerden, ilk ikisi en önemli olan altı raporlama gerekliliği damıttık: koşullama protokolünü belirtin ve eşleştirilmiş olanları yeğleyin; koşulsuz transfer oranlarını mimari karşılaştırması olarak sunmayın. Çalışmamız tek veri kümesini ve tek model çiftini kapsamaktadır; dolayısıyla protokol etkisinin diğer kurulumlardaki boyutu kalibre edilmeyi beklemektedir. Etkinin arkasındaki mekanizma --- hedefler arasındaki eşitsiz temiz doğruluk farkı --- ise kurulumumuza özgü değildir. Protokol karşılaştırmasının yeniden üretilip genişletilebilmesi için tam boru hattını, örnek bazında kayıtları ve analiz betiklerini yayımlıyoruz.
